\documentclass[10pt]{article}
\usepackage[margin=1in]{geometry}
\usepackage{parskip}
\usepackage{xcolor}
\usepackage{hyperref}
% \usepackage{fontspec} 
% remove/replace the above with [T1]{fontenc} if not using XeLaTeX/LuaLaTeX
\usepackage[T1]{fontenc}
\newcommand{\tcolor}{red}
\usepackage{xcolor}
\usepackage{graphicx}
\usepackage{marginnote}
\newcommand{\TODO}[1]{%
  \marginnote{\color{red}\textbf{TODO}}% Put 'TODO' in margin
  {\color{red}#1}% Put the associated text in red in the body
}

% All editable fields (recipient, position, blanks, etc.) live in macros.tex
\input{macros}

\setlength{\parindent}{0pt}
\pagestyle{empty}

\begin{document}

% ---------- Header ----------
\begin{flushright}
\textbf{\ApplicantName}\\ 
% \ApplicantTitle \\
\ApplicantInstitutionNameOne\\
\ApplicantInstitutionNameTwo\\
\ApplicantInstitutionAddressLine\\ \ApplicantInstitutionCityLine\\
\end{flushright}

\RecipientNameOne \hfill \LetterDate \\
\RecipientInstitution \\
% \RecipientDept\\
% \RecipientTitle\\
% \RecipientAddress

% \vspace{1em}
Dear \RecipientNameTwo,

\vspace{0.25em}

I am writing to express my strong interest in the \PositionTitle. 
I am a \ApplicantTitle\ at the \ApplicantInstitutionNameTwo\ in the research group of Prof. Marcelle Soares-Santos, expecting to finish my PhD no later than August 2027.
% my advisor moved from the University of Michigan-Ann Arbor to the University of Zürich, propelling me to 
During my PhD, I have cultivated a unique, independent research program focused on constraining cosmology using electromagnetic and gravitational wave data, developing extensive expertise in astronomical instrumentation, gravitational-wave cosmology, and multi-messenger astrophysics. I believe that my previous work, my expertise in these research areas, and my experience and continued engagement in outreach and science education will \HostGroupDescriptor.

% \textcolor{red}{[Something about the move to Zurich, that was actually to Chile]}
Shortly after beginning my studies in Zurich, I moved to La Serena, Chile to support the Vera C. Rubin Observatory through the final stages of construction and early science operations.
While in Chile, I led the commissioning of the LSST-Camera focal plane array, a SLAC-built instrument. 
\textcolor{\tcolor}{I am eager to apply this expertise to building a new astronomical instrument, designed to capitalize on LSST data, and retain the technical expertise to respond to LSST-Camera anomalies as they arise over the ten-year survey.}

% testing robotic positioners for the DESI instrument and future Stage-V spectroscopic surveys, and in a previous role, delivering the fiber run and spectrograph cameras for the LLAMAS spectrograph for the Magellan Baade telescope. 

In parallel to commissioning the LSST-Camera, I coordinate several observational programs, with a focus on constraining cosmological parameters using observations from optical telescopes and gravitational wave detectors. This includes leading Target-of-Opportunity searches with Vera C. Rubin Observatory and the Dark Energy Camera using the LSST Science pipelines to process obtained data, and augmenting photometric searches using spectroscopy from the Prime Focus Spectrograph. 
% \textcolor{red}{Something here that connects to the host institution.} 
\textcolor{\tcolor}{Retaining this expertise at Stanford, the epicenter of Rubin Observatory operations, will sustain the aforementioned observational programs under a common processing framework. This will help to characterize electromagnetic counterparts of gravitational wave events, and leverage an emerging cosmic probe to gain a deeper understanding of dark energy.} 
%, and the LIGO-Virgo-Kagra gravitational wave observatory. 
% In addition to maximizing data from existing facilities, I am preparing analysis techniques and observing strategies for future observatories, focusing on Stage-V spectroscopic facilities, proposed spectroscopic facilities, and Rubin Observatory to detect the electromagnetic signatures, and forecasted gravitational wave data from Cosmic Explorer, the Einstein Telescope, and the Laser Interferometer Space Antenna. 

% Through these activities, I have become an active member of the leading ground based electromagnetic observatories 
% (Legacy Survey of Space and Time Dark Energy Science Collaboration, Rubin Observatory Project) 
% and gravitational wave observatories
% (LIGO-Virgo-Kagra, the Laser Interferometer Space Antenna).
% , and have established a group of collaborators spanning the United States, Chile, France, Italy, Japan, Great Britain, Germany, Switzerland, and Spain.

Throughout my research career, I have prioritized public engagement and outreach as a core complement to my scientific work. 
During the solar eclipse of 2024, I led a group of high school students from my hometown in guided lectures on solar physics, eclipse safety, and observation of the eclipse from the path of totality using a small solar telescope. 
In 2025, I provided interviews to media outlets covering the first observations of Vera C. Rubin Observatory, including video features with the BBC, and written interviews with the Wall Street Journal and Science News.
% Recently, I led an exhibition at the University of Zürich's Science pavilion during the City of Zürich's annual long night of museums. 
These and other outreach and public engagement activities have reinforced that translating complex physics concepts for non-specialist audiences not only broadens the impact of my research, but also sharpens my own understanding of core physics concepts.
\textcolor{\tcolor}{I am committed to extending these broad impacts by mentoring physics PhD students during their quarterly rotations, mentoring summer SURF students at KIPAC, and engaging with other outreach programs in the wider Menlo Park area.}

Enclosed is my \supplementList. 
% Additionally, \ReferenceOne, \ReferenceTwo, and \ReferenceThree\ would be happy to provide references on my behalf. 
Thank you for your consideration. I welcome the opportunity to discuss how my expertise could contribute to the scientific mission of \RecipientInstitution.

\vspace{0.25em}

All the Best,

\includegraphics[width=0.3\textwidth]{coverLetter/signature.pdf}

\ApplicantName
\end{document}
